%
% 4-homologie.tex
%
% (c) 2026 Prof Dr Andreas Müller
%
\section{Homologie
\label{buch:integration:section:homologie}}
Der Integralsatz von Cauchy besagt, dass ein Deformation des Weges
eines Wegintegrals sich nicht auf den Wert des Integrals auswirkt,
solange die Deformation keine Singulartität der Funktion trifft.
\input{chapters/030-integration/fig/fig-paar.tex}%
Abbildung~\ref{buch:integration:homologie:fig:paar}
zeigt die Deformation des Pfades $\gamma$, der die beiden ``Löcher''
im Definitionsbereich um die Foki $F_1$ und $F_2$ der Lemniskate
einmal umläuft, in zwei Wege $\gamma_+$ und $\gamma_-$, die $F_1$
und $F_2$ separat umlaufen, verbunden über zwei Wegstücke entlang
der reellen Achse, die in entgegengesetzter Richtung durchlaufen
werden und sich im Integral daher wegheben.
Es folgt, dass
\[
\int_\gamma f(z)\,dz
=
\int_{\gamma_+} f(z)\,dz
+
\int_{\gamma_-} f(z)\,dz.
\]
Wir müssen daher die Idee eines Weges wie $\gamma$ zu einer 
Vereinigung von Wegen wie $\{\gamma_+,\gamma_-\}$ verallgemeinern.

\begin{definition}[Kette]
\label{buch:integration:homologie:def:kette}
Eine \emph{Kette} $C$ in $U$ ist eine formale Summe
\index{Kette}%
\begin{equation}
C
=
c_1\gamma_1 + \dots + c_n\gamma_n
\label{buch:integration:homologie:eqn:def:kette}
\end{equation}
von Wegen $\gamma_k\colon I \to U$ mit $c_i\in\mathbb{C}$.
\end{definition}

Das Wegintegral entlang einer Kette wird durch Linearkombination
gebildet:

\begin{definition}[Wegintegral einer Kette]
Ist $C$ eine Kette der Form
\eqref{buch:integration:homologie:eqn:def:kette}
und $f$ eine in $U$ holomorphe Funktion, dann ist das 
\emph{Wegintegral entlang $C$} definiert durch
\index{Wegintegral entlang einer Kette}%
\[
\int_C f(z)\,dz
=
c_1
\oint_{\gamma_1}f(z)\,dz
+
\ldots
+
c_n
\oint_{\gamma_n}f(z)\,dz
=
\sum_{k=1}^n
c_k \oint_{\gamma_k} f(z)\,dz.
\]
\end{definition}

Abbildung~\ref{buch:integration:homologie:fig:paar}
zeigt eine Situation einer Kette $C=\gamma_++\gamma_-$, deren
Wegintegral
\[
\int_C f(z)\,dz
=
\oint_{\gamma_+} f(z)\,dz
+
\oint_{\gamma_-} f(z)\,dz
=
\oint_{\gamma} f(z)\,dz
\]
mit dem Wegintegral entlang der Kurve $\gamma$ übereinstimmt.
Soweit es um Wegintegrale entlang einer Kette geht, gibt es 
zwischen $\gamma$ und $C$ keinen Unterschied.

Zwischen $\gamma$ und nur der Komponente $\gamma_+$ gibt es
ein wesentlichen Unterschied.
Die Funktion $f(z) = 1/(z-F_1)$ hat eine Polstelle im Punkt $F_1$.
Die Wegintegrale sind
\begin{align*}
\frac{1}{2\pi i}
\oint_\gamma f(z)\,dz
&=
\\
\frac{1}{2\pi i}
\oint_\gamma \frac{1}{z-F_1}\,dz
&=
\\
\operatorname{ind}_\gamma(F_1)
&=
1
\\
&\ne
\frac{1}{2\pi i}
\oint_{\gamma_+} f(z)\,dz
\\
&=
\frac{1}{2\pi i}
\oint_{\gamma_+} \frac{1}{z-F_1}\,dz
\\
&=
\operatorname{ind}_{\gamma_+}(F_1)
\\
&=
0
\end{align*}
sind jedoch verschieden.
Der Grund dafür ist, dass $\gamma$ und $\gamma_+$ den Punkt $F_1$
nicht gleich oft umlaufen.
Wir definieren daher ein Äquivalenzrelation von Ketten durch die
Anzahl der Male, die sie Punkte ausserhalb von $U$ umlaufen.

\begin{definition}[Umlaufzahl einer Kette]
Die \emph{Umlaufzahl} einer Kette $C$ um einen Punkt $z_0$ ist
\index{Umlaufzahl}%
\[
\operatorname{ind}_C(z_0)
=
\frac{1}{2\pi i}
\int_C
\frac{1}{z-z_0}
\,dz.
\]
\end{definition}

\begin{definition}[homolog]
Die Ketten $C_1$ und $C_2$ heissen \emph{homolog} in $U$, in Zeichen
$C_1\sim C_2$ oder $C_1\sim_{U}C_2$, wenn
\index{homolog}%
für jeden Punkt $a\in \mathbb{C}\setminus U$
\[
\operatorname{ind}_{C_1}(a)
=
\operatorname{ind}_{C_2}(a)
\]
gilt.
\end{definition}

Die Kette $\gamma_++\gamma_-$ ist also zu $\gamma$ homolog.
Für homologe Ketten ist nicht nur die Umlaufzahl gleich, sondern
jedes Integral.

\begin{satz}
Sind die Ketten $C_1$ und $C_2$ homolog, dann ist
\[
\int_{C_1} f(z)\,dz
=
\int_{C_2} f(z)\,dz.
\]
\end{satz}

Wir betrachten jetzt einen später besonders nützlichen Spezialfall.
\input{chapters/030-integration/fig/fig-merohomolog.tex}%
In Abbildung~\ref{buch:integration:homologie:fig:merohomolog}
ist ein Gebiet $U$ mit einer einfach geschlossenen Kurve
$\gamma$ dargestellt.
Das Komplement von $U$ im Inneren der Kurve $\gamma$ besteht
aus den isolierten Punkten $z_k$.
Jeder dieser Punkte kann so mit der Kurve $\gamma$ verbunden werden,
dass sich die Verbindungen nicht schneiden.
Ausserdem wird jeder von diesen Punkten von einem kleinen Kreis $\gamma_k$
im Gegenuhrzeigersinn umlaufen.
Aus $\gamma$, den im Uhrzeigersinn durchlaufenen Kreisen $-\gamma_k$
und den Verbindungen zu $\gamma$ lässt sich ein neuer Weg zusammensetzen,
der sich in $U$ zusammenziehen lässt.
Das Wegintegral über den gesamten Weg verschwindet daher.
Die Verbindungswege zur Kurve $\gamma$ werden dabei zweimal in
entgegengesetzter Richtung durchlaufen und tragen daher nicht zum
Wegintegral bei.
Das Wegintegral einer in $U$ holomorphen Funktion über $\gamma$
ist daher gleich der Summe der Wegintegrale der Funktion
über die Kreise $\gamma_k$.

\begin{satz}
\label{buch:integration:homologie:satz:pfadpunkte}
Ist $\gamma$ eine einfach geschlossene Kurve in $U$ und besteht das Komplement
von $U$ im Inneren von $\gamma$ nur aus endlich vielen isolierten Punkten
$z_k$, $1\le k\le n$, dann gibt es einen Radius $r$ derart, dass $\gamma$
homolog zu einem Zyklus 
\[
C
=
\sum_{k=1}^n \gamma_k,
\]
wobei jedes $\gamma_k$ ein Kreis vom Radius $r$ um $z_k$ ist.
\end{satz}

\begin{proof}
Es muss nur der Radius $r$ berechnet werden.
Nimmt man
\[
r
=
\frac13
\min \{ |z_k-z_l|\mid 1\le k,l\le n\},
\]
dann können sich die Kreise um die Punkte $z_k$ mit Radius $r$ nicht
schneiden.
Für jeden Punkt $z_k$ ist dann die Umlaufzahl 
\[
\operatorname{ind}_{\gamma_k}(z_l)
=
\delta_{kl}
=
\begin{cases}
1&\qquad\text{für $k=l$}\\
0&\qquad\text{für $k\ne l$.}
\end{cases}
\]
Für den ganzen Zyklus $C$ gilt dann
\[
\operatorname{ind}_C(z_l)
=
\sum_{k=1}^n
\delta_{kl}
=
1
=
\operatorname{ind}_\gamma{z_l}.
\]
Dies zeigt, dass $\gamma$ und $C$ homolog sind.
\end{proof}





